\chapter{Repo et specials}\label{ch:m2:repo-and-specials}

Le lundi 16 septembre 2019, les entreprises américaines ont acquitté leurs impôts trimestriels
et 54~milliards de dollars de bons et d'obligations du Trésor américain nouvellement adjugés ont été réglés.
Ces deux paiements ont ponctionné le cash des comptes que les banques détiennent auprès de la Réserve fédérale~: les réserves ont reculé
d'environ 120~milliards de dollars en deux jours, pour atteindre 1.34~billion de dollars, leur plus bas depuis
2012. Le mardi, le taux des prêts overnight garantis par des \omterm{def:m2:money-markets:bill}{bons du Trésor}, le marché
où chaque dealer finance son inventaire, s'est établi à 5.25~\%. La Réserve
fédérale a fixé pour son taux directeur une fourchette cible de 2.00 à 2.25~\%, et le
taux des fonds fédéraux lui-même a dépassé cette fourchette, pour atteindre 2.30~\%. À 9 h 30 ce matin-là, la New York Fed
a proposé jusqu'à 75~milliards de dollars dans une opération de repo overnight et en a prêté
53~milliards de dollars. Le repo est
le rouage sous-jacent de toute position de cet ouvrage~: une obligation achetée est une
obligation financée, et une obligation vendue à découvert est une obligation empruntée. Ce chapitre suit le cash
et le collatéral à travers lui, et évalue l'unique obligation que tout le monde veut
emprunter.

\section{Le repo vu des deux côtés}

Une pension livrée (repurchase agreement~; livre~1, chapitre~6) est une vente de titres
assortie d'un engagement à les racheter à un prix fixé d'avance à une date ultérieure~: un
échange de collatéral contre du cash, dénoué plus tard, et dont l'écart de prix est
l'intérêt.

\begin{definition}[Taux de repo, reverse repo, repo à terme]\label{def:m2:repo-and-specials:rate}
Le \emph{taux de repo}\index{taux de repo} est le taux d'intérêt simple, en base actual/360
pour le dollar, qui découle de l'écart entre le prix de rachat et le prix de
vente. La partie qui vend le collatéral et emprunte le cash fait un repo~; la
partie qui prête le cash et reçoit le collatéral fait un \emph{reverse
repo}\index{reverse repo}. Un repo overnight arrive à échéance le jour ouvré suivant~;
un \emph{repo à terme}\index{repo à terme} a une échéance fixe plus longue~; un repo ouvert
se reconduit chaque jour jusqu'à ce qu'une des parties y mette fin.
\end{definition}

\begin{example}[Financement du bon à dix ans]\label{ex:m2:repo-and-specials:finance}
Un dealer achète 100~millions de dollars du bon à dix ans décrit au
\cref{ch:m2:government-bonds}, à un prix plein coupon de 100.870911, et le finance
en repo pour 30 jours au taux de 3.90~\%. Avec une décote indicative de 2~\%, le prêteur
paie $100\,000\,000 \times 1.00870911 \times 0.98 = \text{dollars}~98\,853\,492$ et
reçoit 321\,274 d'intérêts à l'échéance~; le dealer finance les 2~\% restants
avec ses fonds propres. Sur les 30 jours, le bon accumule 346\,467~dollars de coupon.
\end{example}

\begin{proposition}[Carry d'une obligation financée]\label{prop:m2:repo-and-specials:carry}
Une position longue de valeur nominale $N$ dans une obligation de coupon $c$ et de prix plein coupon $P$,
financée intégralement au \omterm{def:m2:repo-and-specials:rate}{taux de repo} $r$ pendant $d$ jours d'une période de coupon de $D$
jours, rapporte, si les rendements ne bougent pas,
\[
\text{carry} \;=\; N\,\frac{c}{f}\,\frac{d}{D} \;-\; N\,\frac{P}{100}\,r\,\frac{d}{360},
\]
plus l'effet (faible) de rappel du prix pied de coupon vers le pair. La hausse de rendement qui
l'annulerait est égale au carry divisé par le DV01 de la position.
\end{proposition}
\begin{proof}
Le premier terme est le coupon couru sur les $d$ jours, que le détenteur
encaisse dans le prix plein coupon~; le second est l'intérêt du repo sur le cash
emprunté. Une variation de rendement $\Delta y$ modifie la valeur de $-\mathrm{DV01}\times
\Delta y$ en points de base.
\end{proof}

Au taux de 3.90~\%, le carry du dealer sur le bon à dix ans est $346\,467 - 327\,830 =
\text{dollars}~18\,637$ sur le mois, et une hausse de 0.23 points de base du rendement suffirait à
l'effacer. Un carry positif est la rémunération de la détention d'une obligation qui rapporte
davantage que son financement~; une courbe sur laquelle les obligations longues rapportent moins que le
cash overnight rend le carry négatif, et détenir des obligations revient à parier sur une baisse des rendements.

\section{Repo tripartite, bilatéral et sponsorisé}

\begin{definition}[Repo tripartite et repo bilatéral]\label{def:m2:repo-and-specials:triparty}
Dans un \emph{repo tripartite}\index{repo tripartite}, une banque de compensation agissant en
qualité d'agent détient le collatéral, l'alloue à partir de l'inventaire de l'emprunteur
selon le calendrier d'éligibilité du prêteur, le valorise et applique les
décotes~; le prêt est consenti contre un panier, jamais contre un titre précis.
Dans un \emph{repo bilatéral}\index{repo bilatéral}, les deux parties se règlent directement,
en livraison contre paiement, et le prêteur de cash peut exiger un titre précis.
\end{definition}

\begin{definition}[Repo sponsorisé]\label{def:m2:repo-and-specials:sponsored}
Dans un \emph{repo sponsorisé}\index{repo sponsorisé}, un dealer membre de
la chambre de compensation des \omterm{def:m2:money-markets:bill}{bons du Trésor} sponsorise un non-membre, généralement un \omterm{def:m2:money-markets:mmf}{fonds monétaire}
qui prête du cash, ou un hedge fund qui l'emprunte, dans la chambre de compensation~: la
transaction est novée auprès de la contrepartie centrale et se compense avec les
autres transactions compensées du dealer.
\end{definition}

\begin{omfigure}
\begin{tikzpicture}[font=\footnotesize,
  box/.style={draw,thick,rounded corners=2pt,align=center,minimum height=0.8cm,minimum width=1.8cm,inner sep=3pt},
  arr/.style={-{Stealth[length=2mm]},thick}]
% bilateral
\node[font=\small\bfseries] at (1.6,2.45) {bilatéral};
\node[box,draw=omDef,fill=omDef!8] (l1) at (0.3,1.2) {prêteur de cash};
\node[box,draw=omThm,fill=omThm!8] (b1) at (2.9,1.2) {dealer};
\draw[arr,omDef] (l1.20) -- node[above=2pt] {cash} (b1.160);
\draw[arr,omThm] (b1.200) -- node[below=2pt,align=center] {titre\\précis} (l1.340);
% tri-party
\node[font=\small\bfseries] at (6.6,2.45) {tripartite};
\node[box,draw=omDef,fill=omDef!8] (l2) at (5.3,1.2) {prêteur de cash};
\node[box,draw=omThm,fill=omThm!8] (b2) at (7.9,1.2) {dealer};
\node[box,draw=omProp,fill=omProp!8] (a2) at (6.6,-0.6) {banque agent~:\\alloue, valorise};
\draw[arr,omDef] (l2) -- node[above=9pt] {cash} (b2);
\draw[arr,omThm] (b2) -- node[right=2pt] {panier} (a2);
\draw[arr,omProp] (a2) -- node[left=2pt,align=center] {collatéral\\éligible} (l2);
% sponsored
\node[font=\small\bfseries] at (11.6,2.45) {sponsorisé};
\node[box,draw=omDef,fill=omDef!8] (l3) at (10.3,1.2) {fonds monétaire};
\node[box,draw=omThm,fill=omThm!8] (b3) at (12.9,1.2) {hedge fund};
\node[box,draw=omProp,fill=omProp!8] (c3) at (11.6,-0.6) {chambre de compensation\\(le dealer sponsorise les deux)};
\draw[arr,<->,omDef] (l3) -- node[left=2pt] {novée} (c3);
\draw[arr,<->,omThm] (b3) -- node[right=2pt] {novée} (c3);
\end{tikzpicture}

\omcaption{Trois façons de réaliser un même repo. En \omterm{def:m2:repo-and-specials:triparty}{repo bilatéral}, le prêteur de cash peut
demander un titre précis. Dans un \omterm{def:m2:repo-and-specials:triparty}{repo tripartite}, une banque agent affecte du collatéral
prélevé dans l'inventaire du dealer pour s'adapter à l'échéancier du prêteur. En \omterm{def:m2:repo-and-specials:sponsored}{repo sponsorisé}, un
dealer fait entrer les deux côtés dans la chambre de compensation, qui se porte contrepartie de chacun d'eux,
et les deux transactions du dealer se compensent.}\label{fig:m2:repo-and-specials:structures}
\end{omfigure}

Tout est dans la compensation. Un dealer qui emprunte du cash auprès d'un \omterm{def:m2:money-markets:mmf}{fonds monétaire} dans un repo
et le prête à un hedge fund dans un autre repo détient deux transactions à son bilan
si les deux sont bilatéraux, et, si les deux sont compensés par une même
contrepartie centrale, une position nette proche de zéro. Le bilan d'un dealer
est coûteux, surtout aux dates où il est mesuré (voir au \cref{ch:m2:money-markets}),
la compensation sponsorisée s'est imposée comme un moyen majeur pour les dealers d'intermédier
le repo à grande échelle (voir la \cref{fig:m2:repo-and-specials:sponsored}), et l'obligation
de compensation de la \cref{dat:m2:the-treasury-market:clearing} fera passer beaucoup plus de repo
par la chambre de compensation.

\begin{omfigure}
\begin{tikzpicture}
\begin{axis}[width=11cm,height=4.8cm,
  ylabel={billions de USD},
  tick label style={font=\small},label style={font=\small},
  xmin=2020,xmax=2027,ymin=0,ymax=1.8,grid=major,grid style={gray!25},
  xtick={2020,2021,2022,2023,2024,2025,2026,2027},
  xticklabel style={/pgf/number format/1000 sep=},
  legend columns=2,legend style={font=\footnotesize,at={(0.5,-0.22)},anchor=north,draw=none,fill=none,/tikz/every even column/.append style={column sep=8pt}},
  table/col sep=comma]
\addplot[omDef,very thick] table[x=t,y=reverse]{figdata/markets-2/05-repo-and-specials/sponsored.csv};
\addplot[omThm,very thick] table[x=t,y=repo]{figdata/markets-2/05-repo-and-specials/sponsored.csv};
\legend{reverse repo sponsorisé (cash prêté par les membres sponsorisés),repo sponsorisé (cash emprunté)}
\end{axis}
\end{tikzpicture}

\omcaption{Volumes du \omterm{def:m2:repo-and-specials:sponsored}{repo sponsorisé} FICC à chaque fin de mois, de mars 2020 à août
2026. Les \omterm{def:m2:money-markets:mmf}{fonds monétaires} prêtent via le reverse \omterm{def:m2:repo-and-specials:sponsored}{repo sponsorisé}, les hedge funds empruntent
via le \omterm{def:m2:repo-and-specials:sponsored}{repo sponsorisé}~; les deux ont été multipliés par quatre à cinq, avec des pics aux fins d'année,
où une transaction compensée vaut le plus pour un dealer. Données~: Office of Financial
Research, Hedge Fund Monitor (données FICC).}\label{fig:m2:repo-and-specials:sponsored}
\end{omfigure}

\begin{dated}{2026-09}{Repo sponsorisé FICC}\label{dat:m2:repo-and-specials:sponsored}
Le \omterm{def:m2:repo-and-specials:sponsored}{repo sponsorisé} plus le reverse \omterm{def:m2:repo-and-specials:sponsored}{repo sponsorisé} ont atteint 2.96~billions de dollars le 31
décembre 2025 (1.38~billion de repo et 1.58~billion de \omterm{def:m2:repo-and-specials:rate}{reverse repo}), et
s'élevait à 2.33~billions de dollars le 21 août 2026, la dernière date publiée.
\end{dated}

\section{Specials et specialness}

La plupart des repos sont du \emph{collatéral général} (livre~1, chapitre~16)~: le prêteur de cash
accepte tout titre d'une même classe, et le taux est le prix du cash.
Lorsqu'un acteur a besoin d'un titre particulier, pour le livrer en contrepartie d'une vente à découvert
ou dans un contrat à terme, le taux est aussi le prix de ce titre.

\begin{definition}[Taux de repo special et specialness]\label{def:m2:repo-and-specials:special}
Le \emph{taux de repo special}\index{taux de repo special} d'un titre est le
\omterm{def:m2:repo-and-specials:rate}{taux de repo} auquel les prêteurs de cash prêtent contre ce titre précis. Sa
\emph{specialness}\index{specialness} est le taux de collatéral général moins le
taux special~: ce qu'un prêteur de cash abandonne en intérêts pour obtenir le titre.
\end{definition}

La \omterm{def:m2:repo-and-specials:special}{specialness} est versée au détenteur du titre~: un dealer qui détient le
on-the-run à dix ans peut emprunter du cash contre lui à meilleur coût que contre n'importe quel
autre obligation. Le on-the-run est l'émission la plus vendue à découvert, car les traders
s'en servent pour se couvrir, et c'est aussi celle qui est le plus souvent special. À l'approche
de l'adjudication de son successeur, les ventes à découvert se reportent sur la nouvelle émission, et sa \omterm{def:m2:repo-and-specials:special}{specialness} s'estompe
(voir la \cref{fig:m2:repo-and-specials:special}).

\begin{proposition}[Ce que vaut la specialness]\label{prop:m2:repo-and-specials:value}
Un détenteur qui se finance à une \omterm{def:m2:repo-and-specials:special}{specialness} $s_k$ pendant $d_k$ jours sur des périodes
successives économise, pour 100 de prix plein coupon $P$,
\[
V \;=\; P\,\sum_k \frac{d_k\,s_k}{360}
\]
points de prix. Le titre peut donc se négocier d'autant plus cher qu'un titre par ailleurs
identique, ou $V/\mathrm{DV01}$ points de base de rendement en moins,
avant que sa détention cesse d'être rémunératrice.
\end{proposition}
\begin{proof}
Chaque jour, le détenteur met le titre en repo à $\mathrm{GC} - s_k$ au lieu de GC,
ce qui économise $P\,s_k/360$ pour 100. Sommé sur la durée attendue de la \omterm{def:m2:repo-and-specials:special}{specialness}, cela donne le
gain actualisé, au premier ordre~; un investisseur indifférent entre les deux
titres le paie d'avance dans le prix.
\end{proof}

\begin{omfigure}
\begin{tikzpicture}
\begin{axis}[width=11cm,height=4.5cm,
  xlabel={jours après l'émission},ylabel={specialness (pb)},
  tick label style={font=\small},label style={font=\small},
  xmin=0,xmax=90,ymin=0,ymax=55,grid=major,grid style={gray!25},
  table/col sep=comma]
\addplot[omDef,very thick,const plot] table[x=day,y=bp]{figdata/markets-2/05-repo-and-specials/special.csv};
\node[font=\footnotesize,anchor=south west] at (axis cs:2,46) {nouvelle émission, très vendue à découvert};
\node[font=\footnotesize,anchor=south east] at (axis cs:88,11) {l'adjudication suivante approche};
\end{axis}
\end{tikzpicture}

\omcaption{Trajectoire illustrative de la \omterm{def:m2:repo-and-specials:special}{specialness} d'une nouvelle obligation à dix ans sur les
trois mois précédant l'adjudication de son successeur~: 45 points de base le premier mois,
25 le deuxième, 10 le troisième. L'aire sous la marche est la valeur de
la \cref{prop:m2:repo-and-specials:value}. Les niveaux sont illustratifs. Données~: le
problème du week-end du chapitre.}\label{fig:m2:repo-and-specials:special}
\end{omfigure}

\section{Fails et pénalité de suspens}

Un vendeur à découvert qui ne peut pas emprunter un titre qu'il a vendu manque sa livraison. Jusqu'en
2009, la convention était simplement de reporter la livraison, sans pénalité et à
un prix inchangé~: le défaut ne coûtait que les intérêts sur le produit de la vente, que
l'acheteur ne règle qu'à la livraison. Quand les taux courts se sont effondrés après
l'insolvabilité de Lehman Brothers en septembre 2008, ce coût s'est effondré avec eux
et les fails ont atteint un volume extraordinaire. Le comité des pratiques du marché
a alors instauré une \omterm{def:m2:repo-and-specials:fails}{pénalité de suspens}, à partir de mai 2009.

\begin{definition}[Pénalité de suspens]\label{def:m2:repo-and-specials:fails}
La \emph{pénalité de suspens}\index{pénalité de suspens} d'une transaction sur \omterm{def:m2:money-markets:bill}{bons du Trésor} américain
qui ne se dénoue pas est, pour chaque jour de suspens,
\[
C \;=\; \frac{1}{360}\times 0.01 \times \max(3 - R,\,0) \times P,
\]
où $P$ désigne le produit de la vente et $R$, exprimé en pourcentage, le taux de référence (la borne
inférieure de la fourchette cible de la Federal Reserve) du jour ouvrable précédent. Elle
est versée par le vendeur défaillant à l'acheteur.
\end{definition}

\begin{proposition}[Un plancher pour les taux de repo special]\label{prop:m2:repo-and-specials:floor}
Avec une \omterm{def:m2:repo-and-specials:fails}{pénalité de suspens} de $(3\% - R)^+$ par an, aucun vendeur à découvert n'emprunte un titre à un
\omterm{def:m2:repo-and-specials:special}{taux de repo special} inférieur à $-(3\% - R)^+$~: plutôt que d'accepter moins, il manque sa livraison. Lorsque
$R \ge 3\%$, le plancher est nul~; lorsque $R < 3\%$, les \omterm{def:m2:repo-and-specials:special}{taux de repo special} peuvent être négatifs,
jusqu'à $R - 3\%$.
\end{proposition}
\begin{proof}
Pour livrer, le vendeur à découvert prête des liquidités $C$ au \omterm{def:m2:repo-and-specials:special}{taux de repo special} $s$ et reçoit le
titre, puis le livre et perçoit environ $C$~: il gagne $s$ sur $C$. En revanche, s'il manque sa livraison,
il n'est pas payé, ne gagne donc rien, et paie $(3\% - R)^+$ sur $C$. Il
emprunte le titre si et seulement si $s > -(3\% - R)^+$.
\end{proof}

\begin{dated}{2026-09}{La pénalité de suspens aujourd'hui}\label{dat:m2:repo-and-specials:fails}
La pénalité s'applique aux transactions sur \omterm{def:m2:money-markets:bill}{bons du Trésor} américain conclues depuis le 1\up{er} mai 2009.
La borne inférieure de la fourchette cible est fixée à 3.75~\% depuis le 17 septembre 2026,
$\max(3 - 3.75, 0) = 0$~: manquer la livraison d'un \omterm{def:m2:money-markets:bill}{bon du Trésor} américain n'entraîne actuellement
aucune pénalité, et les \omterm{def:m2:repo-and-specials:special}{taux de repo special} sont minorés à zéro par la proposition ci-dessus.
La pénalité reprend effet si la borne inférieure de la fourchette passe sous 3~\%.
\end{dated}

\section{Septembre 2019}

Les réserves baissaient depuis que la Federal Reserve avait commencé à réduire son
bilan, et elles ont continué à baisser après la fin de ce programme en août 2019~;
dans le même temps, l'émission croissante de Treasuries avait gonflé les inventaires des dealers
de titres à financer. Le 16 septembre, la date fiscale et le règlement des coupons ont retiré
encore 120 milliards de dollars de réserves du système en deux jours. Le \omterm{def:m2:repo-and-specials:rate}{taux
de repo} au jour le jour sur Treasuries a atteint 5.25~\% le 17 septembre
(\cref{fig:m2:repo-and-specials:sept2019}), et le taux des fonds fédéraux est sorti
de sa fourchette cible. L'opération de repo de la Fed de New York à 9~h~30 ce matin-là
a ajouté 53~milliards de dollars de réserves~; le Comité a abaissé sa fourchette le 18, et
à partir de mi-octobre, la Federal Reserve a racheté des \omterm{def:m2:money-markets:bill}{bons du Trésor} à court terme, à raison d'environ 60~milliards de dollars
par mois au départ, afin de maintenir les réserves au moins à leur niveau du début septembre.

\begin{omfigure}
\begin{tikzpicture}
\begin{axis}[width=11cm,height=5cm,
  xlabel={Septembre--octobre 2019},ylabel={taux (\%)},
  tick label style={font=\small},label style={font=\small},
  xmin=0,xmax=42,ymin=1.5,ymax=5.6,grid=major,grid style={gray!25},
  xtick={2,9,16,17,23,30,37},xticklabels={2 sept.,9,16,,23,30 sept.,7 oct.},
  legend columns=3,legend style={font=\footnotesize,at={(0.5,-0.3)},anchor=north,draw=none,fill=none,/tikz/every even column/.append style={column sep=8pt}},
  table/col sep=comma]
\addplot[omThm,very thick,mark=*,mark size=1.2pt] table[x=t,y=sofr]{figdata/markets-2/05-repo-and-specials/sept2019.csv};
\addplot[omDef,very thick] table[x=t,y=effr]{figdata/markets-2/05-repo-and-specials/sept2019.csv};
\addplot[gray,thick,dashed,const plot] table[x=t,y=upper]{figdata/markets-2/05-repo-and-specials/sept2019.csv};
\addplot[gray,thick,dashed,const plot,forget plot] table[x=t,y=lower]{figdata/markets-2/05-repo-and-specials/sept2019.csv};
\legend{SOFR (repo sur Treasuries),taux effectif des fonds fédéraux,fourchette cible}
\node[font=\footnotesize,anchor=west] at (axis cs:18,4.9) {17 septembre~: 5.25~\%};
\end{axis}
\end{tikzpicture}

\omcaption{Taux au jour le jour autour du 17 septembre 2019. Le \omterm{def:m2:repo-and-specials:rate}{taux de repo} garanti
a bondi bien au-dessus de la fourchette du taux directeur~; le taux des fonds fédéraux, un marché
bancaire non garanti, a franchi son plafond de cinq points de base. Après les opérations de la Federal Reserve
et sa baisse du 18, les taux sont revenus dans la nouvelle fourchette,
sauf pour une pointe de fin de trimestre le 30 septembre. Données~: Federal Reserve Bank of
New York et Board of Governors, via FRED.}\label{fig:m2:repo-and-specials:sept2019}
\end{omfigure}

Cet épisode a montré qu'un \omterm{def:m2:central-banks-and-the-short-rate:corridor}{système de plancher} (\cref{ch:m2:central-banks-and-the-short-rate})
a besoin que ses réserves soient abondantes dans chaque banque qui doit les prêter, et non seulement
en moyenne. En juillet 2021, le Comité a créé une facilité de repo permanente, soit une opération
de repo au jour le jour sur Treasuries, à un taux minimum fixe, qui
impose désormais un plafond au \omterm{def:m2:repo-and-specials:rate}{taux de repo} pour les contreparties éligibles.

\section{Tutoriel~: un livre de repo}\label{tut:m2:repo-and-specials:book}

\begin{tutorial}
\textbf{Objectif.} Comptabiliser le financement de l'\cref{ex:m2:repo-and-specials:finance},
calculer son carry et son appel de marge après une baisse de prix, et valoriser le
parcours de \omterm{def:m2:repo-and-specials:special}{specialness} de la \cref{fig:m2:repo-and-specials:special}. \textbf{État
final~:} les chiffres de l'exemple et du problème du week-end.
\begin{tutsteps}
\item \textbf{La transaction}~: les liquidités sont égales à la valeur du collatéral diminuée de la décote~; les intérêts
courent en actual/360.
\omcode{code/firm/repo/firm_repo.py}{11}{32}{Une transaction de repo~: liquidités, intérêts, prix de rachat.}\label{lst:m2:repo-and-specials:trade}
\item \textbf{Marge.} Si le collatéral perd de la valeur, l'emprunteur de liquidités doit
déposer davantage de collatéral pour rétablir la décote.
\omcode{code/firm/repo/firm_repo.py}{39}{45}{L'appel de marge sur un repo après une variation de prix.}\label{lst:m2:repo-and-specials:margin}
Une baisse de 1.5 point sur 100~millions de dollars appelle 1.5~million de dollars le jour même.
\item \textbf{\omterm{def:m2:repo-and-specials:special}{Specialness} et suspens.}
\omcode{code/firm/repo/firm_repo.py}{52}{65}{La valeur de la specialness, la pénalité de suspens et le plancher qu'elle impose aux taux de repo special.}\label{lst:m2:repo-and-specials:special}
Vérifier l'exemple du comité des pratiques~: 5\,555.56~dollars par jour sur 100~millions
de dollars lorsque le taux de référence est de 1~\%.
\item \textbf{Exécuter} \texttt{repo\_demo.financed\_note()} et
\texttt{repo\_demo.special\_value()}~; \texttt{fig\_repo.py} produit les trois
graphiques.
\end{tutsteps}
\textbf{Pour aller plus loin.} Financer la note à 4.30~\% et trouver le
carry négatif~; puis faire doubler la \omterm{def:m2:repo-and-specials:special}{specialness} du premier mois et
exprimer la nouvelle valeur en trente-deuxièmes.
\end{tutorial}

\section{Construction~: le livre de repo}\label{bld:m2:repo-and-specials:book}

\begin{build}
\bfield{Rôle} Chaque position que la firme miniature détient en obligations est
financée ou empruntée~; le livre de repo enregistre les transactions, en calcule les intérêts, émet les appels de marge
et y répond, et indique aux traders ce que coûte la détention de leurs obligations et ce que
leurs ventes à découvert coûtent à emprunter. \bfield{Interface} \texttt{RepoTrade(side, collateral, face, dirty\_price,
haircut, rate, start, end)} avec \texttt{cash}, \texttt{interest(until)},
\texttt{repurchase\_price()}~; \texttt{margin\_call(trade, new\_dirty\_price,
accrued\_to)}~; \texttt{specialness(gc, special)}~;
\texttt{value\_of\_specialness(dirty, path)}~; \texttt{fails\_charge(proceeds,
R, days)}~; \texttt{special\_rate\_floor(R)}~; \texttt{carry(\ldots)}.
\bfield{Règles} Actual/360~; la décote est retenue sur les liquidités~; la marge rétablit
la décote sur les liquidités plus le coupon couru~; la \omterm{def:m2:repo-and-specials:fails}{pénalité de suspens} selon la formule du comité des
pratiques, avec le taux de référence du jour ouvrable précédent.
\bfield{Tests d'acceptation} \texttt{code/firm/repo/tests/}~: liquidités, intérêts et
prix de rachat~; appels de marge dans les deux sens~; l'exemple travaillé du comité des
pratiques au centime près~; le plancher des \omterm{def:m2:repo-and-specials:special}{taux de repo special} au-dessus et au-dessous de 3~\%~; une obligation
dont le coupon égale le financement a un carry nul.
\bfield{Pour aller plus loin} Repos ouverts réévalués chaque jour à partir d'un flux de taux~; substitution de
collatéral~; un calendrier d'éligibilité au \omterm{def:m2:repo-and-specials:triparty}{repo tripartite}~; calcul des intérêts par le
moteur de taux au jour le jour (\cref{bld:m2:central-banks-and-the-short-rate:rfr}).
\end{build}

\begin{omsources}
\item Board of Governors of the Federal Reserve System, FEDS Notes, «~What
Happened in Money Markets in September 2019~?~», 27 février 2020.
\item Treasury Market Practices Group, \emph{Frequently Asked Questions: TMPG
Fails Charges}, 2013.
\item Office of Financial Research, \emph{Hedge Fund Monitor}: FICC sponsored
repo service volumes.
\item K.\,D.\ Garbade, F.\,M.\ Keane, L.\ Logan, A.\ Stokes et J.\ Wolgemuth,
«~The introduction of the TMPG fails charge for U.S. Treasury securities~»,
Federal Reserve Bank of New York \emph{Economic Policy Review}, 2010.
\end{omsources}

\section{Exercices}

\begin{exercise}[$\star$]\label{exo:m2:repo-and-specials:1}
Un fonds prête 50~millions de dollars en repo du vendredi au mardi à 3.90~\%. Calculer
les intérêts qu'il perçoit.
\end{exercise}

\begin{exercise}[$\star$]\label{exo:m2:repo-and-specials:2}
Un nominal de 100~millions de dollars d'un billet dont le prix plein coupon est 101.5 est mis en
repo avec une décote de 2~\%. Calculer le montant de liquidités que reçoit l'emprunteur.
\end{exercise}

\begin{exercise}[$\star$]\label{exo:m2:repo-and-specials:3}
Le collatéral général (GC) se négocie à 3.90~\% et le \omterm{def:m2:repo-and-specials:special}{taux de repo special} du billet à 3.10~\%. Calculer
sa \omterm{def:m2:repo-and-specials:special}{specialness}, ainsi que l'économie de financement que procure chaque jour la détention
de 100~millions de dollars de ce billet.
\end{exercise}

\begin{exercise}[$\star\star$]\label{exo:m2:repo-and-specials:4}
Vérifier le carry de la position à dix ans du dealer décrite dans
l'\cref{ex:m2:repo-and-specials:finance} ainsi que la hausse de rendement qui l'effacerait.
Calculer le carry si le \omterm{def:m2:repo-and-specials:rate}{taux de repo} s'élève à 4.30~\%.
\end{exercise}

\begin{exercise}[$\star\star$]\label{exo:m2:repo-and-specials:5}
Un vendeur ne livre pas pendant cinq jours 100~millions de dollars d'un billet. Calculer la
\omterm{def:m2:repo-and-specials:fails}{pénalité de suspens} qu'il doit lorsque le taux de référence vaut 1~\%, puis lorsqu'il vaut 3.75~\%.
\end{exercise}

\begin{exercise}[$\star\star$]\label{exo:m2:repo-and-specials:6}
En 2021, le taux de référence était de 0~\%. Donner le \omterm{def:m2:repo-and-specials:special}{taux de repo special} le plus bas
qu'un vendeur à découvert accepterait plutôt que de faire défaut, puis celui d'aujourd'hui.
Expliquer ce qui se passait pour les specials avant 2009, lorsque les taux étaient proches de zéro.
\end{exercise}

\begin{exercise}[$\star\star\star$]\label{exo:m2:repo-and-specials:7}
\emph{Programmation.} Avec \texttt{margin\_call}, calculer l'appel de marge sur le repo à 30~jours
du dealer décrit à l'\cref{ex:m2:repo-and-specials:finance} si le prix plein coupon du billet
baisse de 1.5 point le premier jour. Calculer ensuite l'appel de marge après 30~jours si le prix
n'a pas bougé, en tenant compte des intérêts de repo courus.
\end{exercise}

\begin{exercise}[$\star\star\star$]\label{exo:m2:repo-and-specials:8}
\emph{Repérer l'erreur.} « Le taux des fonds fédéraux s'est négocié au-dessus du haut de la
fourchette le 17 septembre 2019, donc le Comité avait resserré la politique. »
Corriger l'affirmation à partir de la \cref{fig:m2:repo-and-specials:sept2019} et du chapitre.
\end{exercise}

\section{Problème~: Le special}

\begin{problem}[{Problème du week-end --- ce que vaut un special}]\label{pb:m2:repo-and-specials:1}
Le billet à dix ans de l'\cref{ex:m2:repo-and-specials:finance} est l'on-the-run.
Un dealer anticipe qu'il se négociera en special de 45~points de base pendant les 30~prochains jours,
de 25~points de base pendant les 30~suivants, et de 10~points de base pendant les 30~derniers jours
précédant la prochaine adjudication, le collatéral général (GC) étant à 3.90~\%. Son prix plein
coupon est de 100.870911 et son DV01 de 0.080444 pour 100.

\medskip\noindent\textbf{Partie I --- Les taux.}
\begin{enumerate}
\item Donner le \omterm{def:m2:repo-and-specials:special}{taux de repo special} dans chacun des trois mois.
\item Quelle économie de financement le dealer réalise-t-il sur 100~millions de dollars le premier
mois~?
\item Et sur les 90~jours~?
\item Pourquoi la \omterm{def:m2:repo-and-specials:special}{specialness} est-elle la plus élevée juste après l'émission~?
\item Qui paie la \omterm{def:m2:repo-and-specials:special}{specialness}, et pourquoi le ferait-il~?
\end{enumerate}

\medskip\noindent\textbf{Partie II --- Le prix.}
\begin{enumerate}[resume]
\item Donner la valeur de la \omterm{def:m2:repo-and-specials:special}{specialness} en points de prix pour 100
(\cref{prop:m2:repo-and-specials:value}).
\item L'exprimer en trente-deuxièmes.
\item L'exprimer en points de base de rendement.
\item Un billet off-the-run de maturité et de coupon presque identiques offre un rendement de 4.210~\%.
En faisant abstraction de la liquidité, à quel rendement l'on-the-run devrait-il se négocier~?
\item Il se négocie à 4.195~\%. Est-il cher ou bon marché sur le seul plan du financement~?
\end{enumerate}

\medskip\noindent\textbf{Partie III --- La transaction.}
\begin{enumerate}[resume]
\item Décrire la transaction qui capture l'écart, jambe par jambe, en précisant
comment chaque jambe est financée.
\item Que rapporte ou que coûte chaque jambe en repo~?
\item Et si la \omterm{def:m2:repo-and-specials:special}{specialness} s'avère deux fois supérieure à la prévision du dealer~?
\item Et si un gros vendeur à découvert provoque un short squeeze sur l'émission, et que son \omterm{def:m2:repo-and-specials:special}{taux
de repo special} tombe à zéro~?
\item Avec un taux de référence à 3.75~\%, jusqu'où le \omterm{def:m2:repo-and-specials:special}{taux de repo special} peut-il descendre~?
\end{enumerate}

\medskip\noindent\textbf{Partie IV --- Jugement.}
\begin{enumerate}[resume]
\item Pourquoi la décote de rendement de l'on-the-run est-elle, sur la plupart des marchés, plus
importante que ne le justifierait sa seule valeur de financement~?
\item Pourquoi un dealer qui détient l'on-the-run et le prête s'intéresse-t-il
à la \omterm{def:m2:repo-and-specials:fails}{pénalité de suspens}~?
\item Comment la décision du Trésor de rouvrir l'émission (en vendre davantage)
modifierait-elle la \omterm{def:m2:repo-and-specials:special}{specialness}~?
\item Énoncer le \emph{résultat clé}~: la valeur de la \omterm{def:m2:repo-and-specials:special}{specialness} attendue en prix, en
trente-deuxièmes et en rendement.
\item En une phrase~: qu'est-ce qu'un special~?
\end{enumerate}
\end{problem}

\section{Questions d'entretien}

\begin{interviewq}[$\star$ \iqroles{trader, banque}]\label{iq:m2:repo-and-specials:1}
Qu'est-ce qu'un repo, et pourquoi est-ce le marché de financement le plus important pour les
traders obligataires~?
\end{interviewq}

\begin{interviewq}[$\star$ \iqroles{trader, chercheur}]\label{iq:m2:repo-and-specials:2}
Que signifie pour une obligation se négocier en special, et qui en profite~?
\end{interviewq}

\begin{interviewq}[$\star\star$ \iqroles{trader, chercheur, banque}]\label{iq:m2:repo-and-specials:3}
Que s'est-il passé sur le marché du repo en septembre 2019, et qu'a fait la Fed~?
\end{interviewq}

\begin{interviewq}[$\star\star$ \iqroles{trader}]\label{iq:m2:repo-and-specials:4}
Vous êtes long d'un billet à dix ans financé en repo. Quel est votre carry, et quel est votre
point mort~?
\end{interviewq}

\begin{interviewq}[$\star\star$ \iqroles{banque, développeur}]\label{iq:m2:repo-and-specials:5}
Pourquoi les dealers préfèrent-ils le \omterm{def:m2:repo-and-specials:sponsored}{repo sponsorisé} et compensé au \omterm{def:m2:repo-and-specials:triparty}{repo bilatéral}~?
\end{interviewq}

\begin{interviewq}[$\star\star\star$ \iqroles{chercheur, trader}]\label{iq:m2:repo-and-specials:6}
Un \omterm{def:m2:repo-and-specials:rate}{taux de repo} peut-il être négatif~? Dériver la borne inférieure d'un \omterm{def:m2:repo-and-specials:special}{taux de repo special}.
\end{interviewq}
